\section{Background}
\label{sec:background}

\paragraph{Setting.}
FLUX.1-dev encodes an $H_{\mathrm{px}}\times W_{\mathrm{px}}$ image with a VAE of stride 8 and packs $2\times2$ latent patches into tokens, so a $1024^2$ image becomes a $64\times64$ grid of tokens with $C=64$ channels. We write a latent as $x\in\R^{C\times H\times W}$ on this token grid and call $64\times64$ the native grid. At $4096^2$ the grid is $256\times256$ and the image carries $N=65{,}536$ tokens instead of $4{,}096$. The per-side scale is $s_a=L_a/64$ for spatial axis $a$ with $L_a$ tokens.

\subsection{RoPE as a Fourier series of the attention logit}
\label{sec:rope-fourier}
FLUX attention heads have dimension $d_k=128$, split into three RoPE axes of 16, 56 and 56 dimensions. The first axis carries a token index that is zero for all image and text tokens in FLUX.1, so its rotations are the identity. The second and third axes encode the row and column of each image token. On a spatial axis with $D=56$ dimensions, RoPE rotates consecutive pairs $(2j,2j{+}1)$, $j=0,\dots,D/2-1$, by the angle $\theta_j p$ at position $p$, with
\begin{equation}
\theta_j = b^{-2j/D},\qquad b=10^4 .
\label{eq:rope-freq}
\end{equation}
Writing each pair as a complex number $\tilde q_j = q_{2j}+i\,q_{2j+1}$, the rotated logit between a query at position $p$ and a key at position $p'$ is
\begin{equation}
\ell(p,p') = \frac{1}{\sqrt{d_k}} \sum_{j} \operatorname{Re}\!\left[\tilde q_j\, \overline{\tilde k_j}\, e^{i\theta_j (p-p')}\right]
          = \frac{1}{\sqrt{d_k}} \sum_{j} |\tilde q_j||\tilde k_j| \cos\!\big(\theta_j \Delta + \phi_j\big),
\label{eq:fourier-logit}
\end{equation}
with $\Delta=p-p'$ and $\phi_j$ the phase of $\tilde q_j\overline{\tilde k_j}$; the sum runs over the pairs of all three axes, and each spatial pair sees the offset along its own axis. For fixed content vectors, the logit is therefore a Fourier series in the offset, with one term per RoPE pair. Pair $j$ has a wavelength of $T_j = 2\pi/\theta_j$ tokens, from $2\pi\approx6.3$ tokens for $j=0$ to about $4.5\times10^4$ tokens for $j=27$ at base $10^4$.

Frequency remapping methods change $\theta_j$. Position interpolation divides every spatial $\theta_j$ by $s_a$. NTK-aware scaling raises the base to $b\,s_a^{D/(D-2)}$, which leaves $\theta_0$ unchanged and compresses the lowest frequency by exactly $s_a$. SEGA uses a stronger base, $b\,s_a^{2D/(D-2)}$~\citep{sega2026}; we call this NTK2. At $s=4$ it gives $b'\approx1.77\times10^5$.

\subsection{Attention as an equalizer}
\label{sec:equalizer}
Multiplying the two components of pair $j$ in both $q$ and $k$ by $m_j>0$ multiplies its term in Eq.~\eqref{eq:fourier-logit} by $g_j=m_j^2$. A vector of per-pair gains is thus a graphic equalizer on the positional Fourier series of the logit. A uniform gain $g$ is an inverse softmax temperature: it sharpens or flattens attention without changing which offsets are preferred. A gain that rises on high-frequency pairs strengthens the terms that discriminate nearby positions and therefore biases attention toward local structure. Because the gain enters as an elementwise scale of $q$ and $k$ after RoPE, it is compatible with fused attention kernels and leaves the softmax scale $1/\sqrt{d_k}$ unchanged.

The softmax normalizer pools over all keys. When the token count grows from $N_{\mathrm{train}}$ to $N$, attention entropy grows roughly with $\log N$ unless logits are scaled up. Jin et al.~\citep{jin2023training} counter this with a logit gain $\log N/\log N_{\mathrm{train}}$, which we write as
\begin{equation}
\tau_N = \frac{\ln N}{\ln N_{\mathrm{train}}},
\label{eq:tau}
\end{equation}
with $N$ counting image tokens. At $4096^2$, $\tau_N = 4/3$.

\subsection{Per-band signal-to-noise ratio of rectified flow}
\label{sec:snr}
FLUX samples a rectified flow~\citep{liu2023flow,lipman2023flow}. With noise level $\sigma\in[0,1]$,
\begin{equation}
x_\sigma = (1-\sigma)\,x_0 + \sigma\,\varepsilon,\qquad \varepsilon\sim\mathcal N(0,I),
\label{eq:flow}
\end{equation}
and the network predicts the velocity $v=\varepsilon-x_0$, from which the clean-image estimate is $\xhat = x_\sigma - \sigma\,v_\theta(x_\sigma)$. Let $\mathcal F$ denote the unnormalized 2D DFT over the token grid, applied per channel. For unit-variance white noise, $\E|\mathcal F\varepsilon_c(k)|^2 = HW$ at every frequency $k$ and channel $c$. Summing over channels, the expected power of $x_\sigma$ at frequency $k$ is
\begin{equation}
\E\,\textstyle\sum_c |\mathcal F x_{\sigma,c}(k)|^2 = (1-\sigma)^2 S(k) + \sigma^2\, C H W,
\label{eq:snr}
\end{equation}
where $S(k)$ is the channel-summed power of $x_0$ and the cross term vanishes in expectation. The per-frequency signal-to-noise ratio is $(1-\sigma)^2S(k)/(\sigma^2CHW)$. Latent spectra fall with frequency~\citep{field1987relations,dieleman2024spectral}, so low frequencies cross unit SNR early in sampling and high frequencies late, which is the coarse-to-fine order of diffusion~\citep{falck2025fourier}. The noise term in Eq.~\eqref{eq:snr} is known exactly at every step. Constancy RoPE uses it as the reference against which it judges whether a band is visible.
